\documentclass[a4paper,11pt,fleqn]{article}%fleqn pour aligner les equations à gauche
\usepackage{Eval}


\begin{document}
    \schemestart
 	\chemfig{H-@{atoh}\charge{60[anchor=180+\chargeangle]=$-$,90=\|,270=\|}{O}}
        +
        \chemfig{H_3C-[1]@{dnl}\charge{45[anchor=180+\chargeangle]=$\delta^+$}{C}(=[2]\charge{-45[anchor=180+\chargeangle]=$\delta^-$,45=\|,135=\|}{O})-[7]\charge{-45[anchor=180+\chargeangle]=$\delta^-$}{O}-[1]CH_2-[7]CH_3}
 	\arrow{->}
 	\chemfig{H-C(-[2]H)(-[6]H)-\charge{45[anchor=180+\chargeangle]=$+$,90[anchor=180+\chargeangle]=\|}{O}(-[6]H)-H}
    \schemestop
  	\chemmove{
 		\draw[shorten <=2pt, shorten >=7pt]
 	(atoh).. controls +(north east:1cm) and +(north west:1cm).. (dnl);}

\end{document}
@{dnl}\charge{-120[anchor=180+\chargeangle]=$\delta^-$,35=\|,125=\|}{O}-[:-35]\charge{-60[anchor=180+\chargeangle]=$\delta^+$}{H}} 

@{atoh}\charge{60[anchor=180+\chargeangle]=$+$}{H}